\documentclass[11pt,a4paper]{article}
\usepackage{curso}
\begin{document}
\portada{Ejercicios --- Lección}{9.4}{Modelo de costo de la difusión}
\section{Ejercicios}
\begin{enumerate}[label=\textbf{E\arabic*.}]
\item (B) Calcule $e_{\text{pos}}$ para $N=7\times10^9$, MFU $=40\,\%$, $e_{\text{FLOP}}=0{,}7$ pJ.
\item (B) Equilibrio $T^\ast$ si la decodificación AR alcanza MFU $=20\,\%$ y la difusión $50\,\%$.
\item (B) TSU: $e_{\text{pos}}$ para 64 celdas, 100 barridos, 2 fJ por celda.
\item (B) ¿Por qué los primeros pasos toleran menos bits?
\item (I) Con lotes grandes, la decodificación AR puede llegar a MFU alto. ¿Cómo cambia $T^\ast$? ¿Qué se pierde (latencia)?
\item (I) Escriba la energía total de un calendario con $b_t$ bits en el paso $t$ si el costo es proporcional a $b_t$.
\item (A) Mida el MFU real de un modelo pequeño en su GPU en decodificación AR y en un paso de difusión, y recalcule $T^\ast$.
\end{enumerate}
\section{Soluciones}
\begin{solucion}
\textbf{E1.} $2\cdot7\times10^9\cdot0{,}7\,\text{pJ}/0{,}4=24{,}5$ mJ.

\textbf{E2.} $T^\ast=0{,}5/0{,}2=2{,}5$.

\textbf{E3.} $64\cdot100\cdot2\ \text{fJ}=12{,}8$ pJ.

\textbf{E4.} La señal es casi puro ruido: los errores de redondeo quedan por debajo del ruido del proceso.

\textbf{E5.} $T^\ast$ baja hacia 1: la ventaja energética de la difusión desaparece, aunque la latencia por usuario de AR crece con el lote.

\textbf{E6.} $E=\sum_te_{\text{pos}}(b_{\max})\,b_t/b_{\max}$.

\textbf{E7.} Depende del hardware; el ejercicio enseña a no usar el pico.
\end{solucion}
\end{document}
